% !TeX root = appendix-of-notes.tex
\section{Appendix}


\subsection{Geometric invariant theory}

Here we provide a brief overview of geometric invariant theory (GIT).
For the propose of this notes, we only consider that the ground field $\kkk$ is algebraically closed and of characteristic zero, typically $\kkk = \bbC$.

\begin{definition}[Categorical and geometric quotients]
	Let $G$ be an algebraic group acting on a variety $X$ via $\sigma: G \times X \to X$.
	A \emph{categorical quotient} of $X$ by $G$ is a variety $Y$ together with a $G$-invariant morphism $\pi: X \to Y$ such that for any other variety $Z$ and any $G$-invariant morphism $f: X \to Z$, there exists a unique morphism $\tilde{f}: Y \to Z$ making the following diagram commute:
	\[
		\begin{tikzcd}
			G \acts X & Y \racts \id \\
			& Z \racts \id
			\arrow[from=1-1, to=1-2, "\pi"]
			\arrow[from=1-2, to=2-2, "\tilde{f}", gray]
			\arrow[from=1-1, to=2-2, "f"]
		\end{tikzcd}
	\]

	Given a categorical quotient $\pi: X \to Y$, we say that $Y$ is a \emph{geometric quotient} of $X$ by $G$ if the following conditions hold:
	\begin{enumerate}
		\item The variety $Y$ is equipped with the quotient topology induced by $\pi$, i.e., a subset $U \subseteq Y$ is open if and only if $\pi^{-1}(U)$ is open in $X$.
		\item The morphism $\pi$ is surjective, and the fibers of $\pi$ are precisely the $G$-orbits in $X$, i.e., for each $y \in Y$, we have $\pi^{-1}(y) = G \cdot x$ for some $x \in X$.
		\item The sheaf of regular functions on $Y$ is isomorphic to the sheaf of $G$-invariant regular functions on $X$, i.e., for any open subset $U \subseteq Y$, we have $\mathcal{O}_Y(U) \cong \mathcal{O}_X(\pi^{-1}(U))^G$.
	\end{enumerate}
\end{definition}

In general, a categorical quotient may not be a geometric quotient, and the existence of quotients is not guaranteed.
\begin{example}
	Let $G = \bbG_m$ act on $X = \bbA^1$ by scaling, i.e., $t \cdot x = tx$ for $t \in \bbG_m$ and $x \in \bbA^1$.
	The categorical quotient of this action is a single point, since any $G$-invariant morphism from $\bbA^1$ to another variety must be constant.
	Easily, we can see that the categorical quotient is not a geometric quotient, since the fibers of the quotient map are not precisely the $G$-orbits in $X$.
\end{example}

On GIT, we are particularly interested in the case where $G$ is a reductive group.

Recall that any affine linear algebraic group $G$ is a linear algebraic group.
For a linear algebraic group $G$ and $g \in G$, we say that $g$ is \emph{semisimple} if it is diagonalizable over the algebraic closure of the base field, and $g$ is \emph{unipotent} if all its eigenvalues are equal to 1.
We have the \emph{Jordan decomposition} of $g$ as $g = g_s g_u$, where $g_s$ is semisimple and $g_u$ is unipotent, and they commute with each other.

On the other hands, $G$ has a unique maximal normal solvable connected subgroup, which is called the \emph{radical} of $G$, denoted by $R(G)$.
All unipotent elements of $R(G)$ form a normal subgroup of $G$, called the \emph{unipotent radical} of $G$, denoted by $R_u(G)$.
The unipotent radical $R_u(G)$ can also be characterized as the unique maximal normal unipotent connected subgroup of $G$.
We say that $G$ is \emph{reductive} if its unipotent radical is trivial, i.e., $R_u(G) = \{e\}$, where $e$ is the identity element of $G$.
This is equivalent to saying that the radical $R(G)$ is a torus, i.e., $R(G) \cong (\bbG_m)^n$ for some $n \geq 0$.

\begin{example}
	The general linear group $\GL_n$ is reductive, since its radical is the center of $\GL_n$, which is isomorphic to $\bbG_m$.

	To see this, let $R = R(\GL_n)$ be the radical of $\GL_n$.
	By Lie-Kolchin theorem, $R$ is conjugate to a subgroup of the group $B$ of upper triangular matrices in $\GL_n$.
	Since $R$ is normal in $\GL_n$, it is contained in the intersection of all conjugates of $B$, which is the group of scalar matrices, i.e., $R \subseteq Z(\GL_n) \cong \bbG_m$.
	Note that we always have $Z(G)^\circ \subseteq R(G)$ for any linear algebraic group $G$, where $Z(G)^\circ$ is the connected component of the identity in the center of $G$.
	Hence $R$ is a torus, and thus $\GL_n$ is reductive.

	By similar argument, one can show that $R(\SL_n) = \{e\}$, and thus $\SL_n$ is also reductive.
\end{example}

For reductive groups, we have the following theorem to show its power.

\begin{theorem}[{ref. \cite[Theorem 1.1, Amplification 1.3]{mumford1994-git}}]
	Let $G$ be a reductive group acting on an affine variety $X$.
	Then the categorical quotient $X \to X / G$ exists.
	And it is a geometric quotient if and only if all $G$-orbits (of closed points) in $X$ are closed.
\end{theorem}

For projective cases, the case is a bit more complicated, and we need to introduce the notion of $G$-linearization and (semi)stable points.

\begin{definition}
	Let $G$ be an algebraic group acting on a variety $X$ and let $\calL$ be a line bundle on $X$.
	A \emph{$G$-linearization} of $\calL$ is an isomorphism $\phi: \sigma^* \calL \to p_2^* \calL$ of line bundles on $G \times X$, where $\sigma: G \times X \to X$ is the action map and $p_2: G \times X \to X$ is the projection onto the second factor.
\end{definition}

Inspiringly, assume that $\calL$ is very ample with an embedding $\iota: X \hookrightarrow \bbP^n$.
Then given a $G$-linearization $\phi: \sigma^* \calL \to p_2^* \calL$, by restricting $\phi$ to the $\{g\} \times X$ for each $g \in G$, we get an isomorphism $\phi_g: \sigma_g^* \calL \to \calL$, where $\sigma_g: X \to X$ is the action of $g$ on $X$.
Hence we induce a (contravariant) linear action of $G$ on $H^0(X, \calL)$, and thus a linear action of $G$ on $\bbP^n$.
In this case, a $G$-linearization of $\calL$ will induce a $G$-equivariant embedding $\iota: X \hookrightarrow \bbP^n$.

\begin{definition}
	Let $G$ be a reductive group acting on a projective variety $X$ and let $\calL$ be a $G$-linearized line bundle on $X$.
	A point $x \in X$ is called \emph{semistable} if there exists a $G$-invariant section $s \in H^0(X, \calL^{\otimes n})^G$ for some $n > 0$ such that $s(x) \neq 0$ and $X_s := \{y \in X \mid s(y) \neq 0\}$ is affine.
	A semistable point $x \in X$ is called \emph{stable} if the action of $G$ on $X_s$ is closed, i.e., the orbit $G \cdot y$ is closed in $X_s$ for all $y \in X_s$.
\end{definition}

The set of semistable points and stable points are open in $X(\kkk)$, and we will also denote the corresponding open subvarieties by $X^{\text{ss}}(\calL)$ and $X^{\text{s}}(\calL)$ respectively.

Then we have the following theorem in projective cases.

\begin{theorem}[{ref. \cite[Theorem 1.10 and comments under Amplification 1.11]{mumford1994-git}}]
	Let $G$ be a reductive group acting on a projective variety $X$ and let $\calL$ be a $G$-linearized line bundle on $X$.
	\begin{enumerate}
		\item there exists a categorical quotient $\pi: X^{\text{ss}}(\calL) \to X^{\text{ss}}(\calL) / G$;
		\item $X^{\text{ss}}(\calL) / G \cong \Proj \left( \bigoplus_{n \geq 0} H^0(X, \calL^{\otimes n})^G \right)$;
		\item the restriction of $\pi$ to $X^{\text{s}}$ is a geometric quotient $\pi: X^{\text{s}}(\calL) \to X^{\text{s}}(\calL) / G$.
	\end{enumerate}
\end{theorem}

The next problem is to determine the (semi)stable points.

\begin{construction}
	Settings as the theorem above.
	Let $l: \bbG_m \to G$ be a homomorphism of algebraic groups, which is called a \emph{$1$-parameter subgroup} of $G$.
	Then we have an action of $\bbG_m$ on $X$ via $l$, and thus we can define the limit point $\lim_{t \to 0} l(t) \cdot x$ for any $x \in X$.
	More exactly, given an $1$-parameter subgroup $l: \bbG_m \to G$, we can define a morphism $\rho_{x, l}: \bbG_m \to X$ by $\rho_{x, l}(t) = l(t) \cdot x$.
	And then by the valuative criterion of properness, we can uniquely extend $\rho_{x, l}$ to a morphism $\tilde{\rho}_{x, l}: \bbA^1 \to X$, and thus we can define the limit point $\lim_{t \to 0} l(t) \cdot x = \tilde{\rho}_{x, l}(0)$.
	This limit point will be denoted by $l(0) \cdot x$.

	For every $\alpha \in \bbG_m$, we can shift the $1$-parameter subgroup $l$ by $\alpha$ to get another $1$-parameter subgroup $l_\alpha: \bbG_m \to G$ defined by $l_\alpha(t) = l(\alpha t)$.
	Note that $l_\alpha(0) \cdot x = l(0) \cdot x$ for all $\alpha \in \bbG_m$ since $\tilde{\rho}_{x, l_\alpha}(t) = \tilde{\rho}_{x, l}(\alpha t)$.
	Hence $l(\alpha) \cdot (l(0) \cdot x) = l(0) \cdot x$ for all $\alpha \in \bbG_m$.

	Note that since $\calL$ is $G$-linearized, we have an induced action of $G$ on the vector bundle $\bbV(\calL)$, and thus we have an action of $\bbG_m$ on the fiber $\calL_{l(0) \cdot x}$.
	This gives a character $\chi: \bbG_m \to \bbG_m$ defined by $l(\alpha) \cdot v = \chi(\alpha) v$ for all $v \in \bbV(\calL)|_{l(0)\cdot x}$ and $\alpha \in \bbG_m$.
	Identifying $\Hom(\bbG_m, \bbG_m) \cong \bbZ$, we can write $\chi(\alpha) = \alpha^r$ for some $r \in \bbZ$.
	Then we set $\mu^\calL(x, l) = r$.
\end{construction}

\begin{remark}
	Note that $\bbV(\calL) = \Spec \Sym^\bullet \calL$ is in the sense of Grothendieck, and the fiber $\bbV(\calL)|_x$ is dual to the fiber $\calL|_x$.
	Hence our definition of $\mu^\calL(x, l)$ coincides with the one in \cite[Definition 2.2]{mumford1994-git} although we have a sign difference.
	\Yang{}
\end{remark}

\begin{theorem}[{ref. \cite[Theorem 2.1]{mumford1994-git}}]
	Let $G$ be a reductive group acting on a projective variety $X$ and let $\calL$ be a $G$-linearized line bundle on $X$.
	Then we have the following criterion for (semi)stability:
	\begin{enumerate}
		\item A point $x \in X$ is semistable if and only if $\mu^\calL(x, l) \geq 0$ for all $1$-parameter subgroups $l: \bbG_m \to G$.
		\item A point $x \in X$ is stable if and only if $\mu^\calL(x, l) > 0$ for all $1$-parameter subgroups $l: \bbG_m \to G$.
	\end{enumerate}
\end{theorem}

If $X = \bbP^n$ and $\calL = \calO_{\bbP^n}(1)$, then we can give a more explicit description of $\mu^\calL(x, l)$.

\begin{proposition}[{ref. \cite[Example 2.3]{mumford1994-git}}]
	Let $G$ be a reductive group act on $X = \bbP^n$ and $\calL = \calO_{\bbP^n}(1)$.
	Then for any $1$-parameter subgroup $l: \bbG_m \to G$, there exists a coordinate system such that $l(t) = \diag(t^{r_0}, t^{r_1}, \ldots, t^{r_n})$ with $\sum_{i=0}^n r_i = 0$.
	Under this coordinate system, for any point $x = [x_0: x_1: \ldots: x_n] \in \bbP^n$, we have
	\[
		\mu^\calL(x, l) = -\min\{r_i \mid x_i \neq 0\}.
	\]
\end{proposition}

\begin{example}
	Let \Yang{}
\end{example}


\subsection{Tools in complex dynamics}

On this subsection, we review some tools in complex dynamics quickly, in particular, the notion of bifurcation measure, bifurcation current, Green current, plurisubharmonic functions and so on.




































